% META: {"id":"1SPE-VARALEA-CO-046","chapitre":"1SPE-VARIABLES-ALEATOIRES","type_objet":"corrige","exercice_ref":"1SPE-VARALEA-EX-046","capacites_codes":["C1","C2","C3","C4","C5"],"sources_inspiration":[],"status":"approved","origin":"generated","status_history":["generated","audited","accepted","approved"],"release_acceptance":"RELEASE_OWNER_BATCH_ACCEPTANCE","acceptance_closure_digest":"sha256:9b3ccf9a81c5520fbb7e03b7b2d3e7bf2057a02834b6d908bf3be5f49f3b553f","acceptance_receipt_digest":"sha256:4fad86f0282e0fa4e0419cca1ad6379ae7af5a280c04a4a90880f90a1c044242"}
% BEGIN-VERIFY
% from sympy import Rational
% probabilities = [Rational(125,216), Rational(75,216), Rational(15,216), Rational(1,216)]
% E_X = sum(k*probabilities[k] for k in range(4)); E_G = 10*E_X - 5
% assert E_G == 0
% E_X2 = sum(k*k*probabilities[k] for k in range(4))
% V_X = E_X2-E_X**2; assert V_X == Rational(5,12)
% probabilities_G = [Rational(125,216), Rational(75,216), Rational(15,216), Rational(1,216)]
% values_G = [-5, 5, 15, 25]
% E_G2 = sum(g*g*p for g,p in zip(values_G, probabilities_G))
% V_G = E_G2-E_G**2; assert V_G == Rational(125,3)
% END-VERIFY
\begin{corrige}{1SPE-VARALEA-EX-046}

\textbf{1.} Les $216$ issues ordonnées sont équiprobables. Il y a respectivement
$125$, $75$, $15$ et $1$ chemins comportant zéro, un, deux et trois résultats
égaux à $6$. La loi de $X$ est donc donnée par la ligne de probabilités du tableau suivant.

\medskip
\textbf{2.}
\[
\begin{array}{|c|c|c|c|c|}
\hline
k & 0 & 1 & 2 & 3 \\
\hline
G = 10k-5 & -5 & 5 & 15 & 25 \\
\hline
P & \frac{125}{216} & \frac{75}{216} & \frac{15}{216} & \frac{1}{216} \\
\hline
\end{array}
\]

\medskip
\textbf{3.} D'après le tableau, $E(X)=\frac12$. Donc
$E(G)=10E(X)-5=10\times\frac12-5=0$. Le jeu est équitable.

\medskip
\textbf{4.} À partir du tableau de loi de $G$,
\[
E(G^2)=25\times\frac{125}{216}+25\times\frac{75}{216}
+225\times\frac{15}{216}+625\times\frac{1}{216}=\frac{125}{3}.
\]
Comme $E(G)=0$, on obtient $V(G)=E(G^2)-[E(G)]^2=\frac{125}{3}
\approx41{,}7$ et $\sigma(G)=\sqrt{\frac{125}{3}}\approx6{,}45$~euros.

\end{corrige}
